\section{Реализация программы}

Реализация программы представлена в листинге \ref{lst:prog}.

\begin{lstlisting}[caption={Реализация программы}, label={lst:prog}]
factorial_tail(0, Acc, Acc) :- !.
factorial_tail(N, Acc, Result) :-
    N > 0,
    NewAcc is Acc * N,
    NewN is N - 1,
    factorial_tail(NewN, NewAcc, Result).

factorial(N, Result) :- factorial_tail(N, 1, Result).

fib_tail(0, Acc1, _, Acc1) :- !.
fib_tail(N, Acc1, Acc2, Result) :-
    N > 0,
    NewAcc2 is Acc1 + Acc2,
    NewN is N - 1,
    fib_tail(NewN, Acc2, NewAcc2, Result).

fibonacci(N, Result) :- fib_tail(N, 0, 1, Result).
\end{lstlisting}

Форма вопроса вида \ref{lst:query1} поиска n!.
\begin{lstlisting}[caption={Форма задания вопроса 1}, label={lst:query1}]
?- factorial(n, Res).
\end{lstlisting}

Форма вопроса вида \ref{lst:query2} поиска n-го числа Фибоначчи.
\begin{lstlisting}[caption={Форма задания вопроса 2}, label={lst:query2}]
?- fibonacci(n, Res).
\end{lstlisting}

\clearpage
\section{Варианты вопросов и ответов на них}

\begin{lstlisting}[caption={Вопрос и ответ №1}, label={lst:q1}]
?- factorial(0, Res).
Res = 1.
\end{lstlisting}

\begin{lstlisting}[caption={Вопрос и ответ №2}, label={lst:q2}]
?- factorial(1, Res).
Res = 1.
\end{lstlisting}

\begin{lstlisting}[caption={Вопрос и ответ №3}, label={lst:q3}]
?- factorial(4, Res).
Res = 24.
\end{lstlisting}

\begin{lstlisting}[caption={Вопрос и ответ №4}, label={lst:q4}]
?- fibonacci(0, Res).
Res = 0.
\end{lstlisting}

\begin{lstlisting}[caption={Вопрос и ответ №5}, label={lst:q5}]
?- fibonacci(1, Res).
Res = 1.
\end{lstlisting}

\begin{lstlisting}[caption={Вопрос и ответ №6}, label={lst:q6}]
?- fibonacci(2, Res).
Res = 1.
\end{lstlisting}

\begin{lstlisting}[caption={Вопрос и ответ №7}, label={lst:q7}]
?- fibonacci(6, Res).
Res = 8.
\end{lstlisting}
